\chapter{מטרות ויעדים}\label{ch:goals}

\section{מטרת על}
\begin{keybox}[מטרת על]
לאפשר לכל מערכת חדשות להבטיח ולהוכיח לקוראיה שכל תמונה שהיא מציגה כתיעוד היא אכן תיעוד, ושכל תמונה סינתטית מוצגת ככזו.
\end{keybox}

\section{יעדים מדידים}
כל יעד מנוסח לפי \en{SMART}: ספציפי, מדיד, בר השגה, רלוונטי ותחום בזמן. העמודה ״שיטת מדידה״ קובעת מראש איך ייבדק היעד, כדי שבסוף הפרויקט התשובה תהיה ״עומד״ או ״לא עומד״, בלי פרשנות.

\begin{table}[!htbp]\centering
\caption{יעדי הפרויקט, מדדי ההצלחה ושיטת המדידה}\label{tab:goals}
\small
\begin{tabularx}{\linewidth}{@{}l Y l Y l@{}}
\toprule
\textbf{מזהה} & \textbf{יעד} & \textbf{מדד הצלחה} & \textbf{שיטת מדידה} & \textbf{מועד} \\
\midrule
\req{G-1} & זיהוי אמין של תמונות סינתטיות & גילוי \num{90\%} לפחות, התרעות שווא \num{2\%} לכל היותר & מבחן על מאגר \en{VL-Wild} (סעיף~\ref{sec:bench}) & חודש \num{12} \\
\req{G-2} & החלטה מהירה בזמן אמת & חציון אוטומטי $\leq$ \num{60} שנ׳; עם אדם $\leq$ \num{15} דק׳ & מדידה מיומני המערכת בפיילוט & חודש \num{15} \\
\req{G-3} & צמצום פרסום תמונות שקריות & ירידה $\geq$ \num{80\%} לעומת קו בסיס & ביקורת מדגמית של \num{2,000} תמונות לפני ואחרי & חודש \num{18} \\
\req{G-4} & שקיפות מלאה לקורא & \num{100\%} מהתמונות שפורסמו מתויגות לפי \en{TRUTH-IMG} & סריקה אוטומטית של עמודי האתר & חודש \num{18} \\
\req{G-5} & אימוץ התקן & \num{5} מערכות פיילוט מוסמכות; \num{20} מתחייבות & חוזים חתומים ודוחות ביקורת & חודש \num{18} \\
\req{G-6} & מוכנות לקנה מידה & קיבולת של \num{2} מיליון תמונות בחודש במבחן עומס; זמינות \num{99.9\%} & מבחן עומס וניטור זמינות & חודש \num{16} \\
\bottomrule
\end{tabularx}
\end{table}

\section{יעדי קנה מידה (אחרי הפרויקט)}
יעדים אלה אינם חלק מקריטריוני הקבלה של הפרויקט, אבל הארכיטקטורה נבחנת מולם כבר עכשיו, כדי שהמוצר לא יידרש לבנייה מחדש כשיגדל:
\begin{itemize}
\item \textbf{36 חודשים:} \num{1,000} מערכות חדשות ב-\num{30} מדינות; \num{50} מיליון תמונות מאומתות בחודש.
\item \textbf{60 חודשים:} \en{TRUTH-IMG} מוכר כתקן תעשייתי על ידי לפחות גוף אתיקה לאומי אחד ורשת פרסום גלובלית אחת.
\end{itemize}

\section{מה הפרויקט אינו מבטיח}
כחלק מהחוזה, נקבע גם מה \textbf{אינו} מובטח: המערכת אינה מבטיחה זיהוי של \num{100\%} מהזיופים, וגם לא החלטה אוטומטית בלי אדם. ציון ״לא ודאי״ הוא תוצאה לגיטימית ומכוונת, שמעבירה את ההחלטה לאדם.
